\appendix
\setcounter{equation}{0}
\renewcommand{\theequation}{A\arabic{equation}}
\renewcommand{\theHequation}{A.\arabic{equation}}

\section{Analytic Occurrence Integrals}\label{app:analytic}

For exponent $q$, define
\begin{equation}
 \mathcal{J}(x_0,x_1;q)=
 \begin{cases}
 \displaystyle\frac{x_1^{q+1}-x_0^{q+1}}{q+1}, & q\neq-1,\\[6pt]
 \ln(x_1/x_0), & q=-1.
 \end{cases}
 \label{eq:primitive}
\end{equation}
With $T_0=3900$ K, $T_b=5117$ K, and $T_1=6300$ K, the Model~1 normalization is
\begin{align}
 C_1^{-1}={}&
 \mathcal{J}(0.5,2.5;\alpha)
 \mathcal{J}(0.2,2.2;\beta)\nonumber\\
 &\times\frac{1}{T_1-T_0}\big[
 10^{-11.839}\mathcal{J}(T_0,T_b;\gamma+3.16)\nonumber\\
 &\hspace{36mm}+10^{-16.769}\mathcal{J}(T_b,T_1;\gamma+4.49)
 \big].
 \label{eq:c1}
\end{align}
Because Equation~(\ref{eq:lambda}) is separable within each temperature branch, $f_k(T)$ is evaluated analytically in $r$ and $I$ using Equation~(\ref{eq:primitive}). The only numerical quadrature in the final Galactic estimator is the radial integration after the row-level sums have been constructed.

\section{Bryson Model~2 Sensitivity}\label{app:model2}

For the Model~2 sensitivity, the source occurrence function is
\begin{equation}
 \lambda_2(r,I,T\mid\bm\theta_2)=F_0C_2r^\alpha I^\beta T^\gamma,
 \label{eq:model2}
\end{equation}
without the fixed geometric factor $g(T)$. Its normalization is
\begin{equation}
 C_2^{-1}=\mathcal{J}(0.5,2.5;\alpha)
 \mathcal{J}(0.2,2.2;\beta)
 \frac{\mathcal{J}(T_0,T_1;\gamma)}{T_1-T_0}.
 \label{eq:c2}
\end{equation}
The adopted \citet{Bryson2021} \texttt{hab2}, Poisson, ``With Uncertainty,'' constant-completeness median vector is
\begin{equation}
 (F_0,\alpha,\beta,\gamma,C_2)
 =(1.13,-1.13,-0.85,+1.19,1.072724\times10^{-5}).
\end{equation}
Evaluated on the same frozen JJ rows and HZ prescription, this gives $\LamEE=3.487759\times10^6$, or $+3.296\%$ relative to the Model~1 constant-completeness reference.

\section{Frozen and Native PARSEC-TAMS Boundaries}\label{app:tams}

Table~\ref{tab:tams_nodes} lists the inherited monotonic segment used by the
canonical 5300--6000 K selector. The 5200 K, $1.15R_\odot$ row is retained
verbatim as source provenance. An independent regeneration uses genuine
solar-composition low-mass phase-7 nodes at 5151.34 K ($0.75\,M_\odot$),
5304.81 K ($0.80\,M_\odot$), and 5390.14 K ($0.85\,M_\odot$), followed by the
native hotter sequence. This curve brackets the full target without the 5200 K
row and produces no integrated selector disagreement. Thus the special anchor
is not numerically responsible for the canonical host count.

\begin{table}[h]
\caption{Frozen TAMS nodes used for G-star interpolation}
\label{tab:tams_nodes}
\centering
\begin{tabular}{rr}
\toprule
$T_{\rm eff}$ (K) & $R_{\rm TAMS}/R_\odot$\\
\midrule
5200.00000 & 1.15000\\
5390.13944 & 1.22926\\
5517.85139 & 1.28542\\
5633.13293 & 1.35053\\
5738.25706 & 1.42375\\
5844.13178 & 1.49188\\
5951.82290 & 1.55332\\
6060.24246 & 1.61155\\
\bottomrule
\end{tabular}
\end{table}

\section{Frozen Verification Checklist}\label{app:checks}

The v4 freeze requires: (1) input and artifact checksums; (2) bitwise legacy
measurement regression and corrected quantile recovery; (3) post-perturbation
source-domain filtering; (4) 400 of 400 converged adaptive chains in each
branch; (5) equal outer-realization weights; (6) whole-realization outer MCSE
and contiguous-block inner MCSE below their declared thresholds; (7)
independent MCMC-seed stability; (8) finite non-negative host weights and
occurrence values; (9) radial-node, unit, and thin/thick closure; (10) correct
temperature, age, component, TAMS, and compact-remnant selection; (11) nested
HZ/target domains and analytic occurrence integrals; (12) native solar TAMS
agreement without the 5200 K anchor; (13) 0.5-to-0.25 kpc radial convergence;
(14) exact DR25 target-support masks; and (15) estimand-aware separation of
posterior scenarios, point alternatives, and categorical failures.

The metallicity-dependent TAMS artifact does not pass this checklist. Its
high-mass/giant contamination and missing low-metallicity temperature coverage
are recorded as a failed, excluded experiment. Likewise, the zero-candidate
DR25 result is a scientific support failure, not a software failure. These
distinctions prevent engineering reproducibility from being mislabeled as
observational validation.

\bibliography{references}
\bibliographystyle{aasjournalv7}
